% =============================================================
%  GENERAL INTRODUCTION
% =============================================================

\chapter*{General Introduction}
\addcontentsline{toc}{chapter}{General Introduction}
\markboth{General Introduction}{}

\setstretch{1.5}

\noindent
The digital revolution has redrawn the boundaries of how knowledge is produced, accessed and exchanged. In a few short years, the act of \textit{learning} has shifted from an exclusively physical experience -- bound by the walls of a classroom -- towards a flexible, connected and increasingly personalized journey, available anywhere and at any time. This silent transformation has placed online training at the very heart of professional growth, lifelong learning and competitiveness in the labor market.

\medskip
\noindent
In Tunisia, training centers and technological institutes are at a turning point. The expectations of learners have evolved : they no longer seek a static catalog of courses, but a smart, interactive and personalized environment, capable of guiding them, tracking their progress and helping them make the right choices. At the same time, training organizations need powerful management tools to handle their users, their offerings, their reservations and their administrative workflows in a fluid and reliable way.

\medskip
\noindent
It is in this context that our end-of-studies project takes place, carried out within \textbf{ADVANCIA Training Center} as part of our Applied Bachelor's Degree in Computer Science at the \textbf{Institut Supérieur des Études Technologiques de Siliana}. The mission entrusted to us was clear in its ambition : to design and develop a \textbf{smart web platform for training management}, named \textbf{ADVANCIA}, capable of bringing together \textit{users}, \textit{trainers}, \textit{administrators} and \textit{super-administrators} within a single, secure and modern ecosystem.

\medskip
\noindent
Our solution is built around four well-defined actors. The \textbf{User} can create an account, browse the available trainings, reserve sessions, manage his profile and follow his personal dashboard. The \textbf{Trainer} animates the courses he is responsible for and follows his learners. The \textbf{Administrator} pilots the platform on a daily basis : managing users, validating reservations, importing and exporting data in Excel or PDF, and exploiting a smart dashboard supported by an integrated chatbot. Finally, the \textbf{Super-Administrator} holds the highest level of authority, supervising every actor and every operation of the system.

\medskip
\noindent
To structure this work in a coherent and progressive manner, the present dissertation is organized into six chapters that reflect the methodological journey we followed.

\begin{itemize}[leftmargin=2.2em,itemsep=0.4em]
  \item \textbf{Chapter 1 -- Project Framework Presentation :} introduces the host company \textsc{ADVANCIA Training Center}, situates the project in its context, identifies the problem and presents the proposed solution along with the chosen methodology.
  \item \textbf{Chapter 2 -- Analysis and Project Planning :} captures the functional and non-functional requirements, identifies the actors, formalizes the global use case diagram, and structures the work into releases and sprints supported by a Gantt schedule.
  \item \textbf{Chapter 3 -- Release 1 :} addresses the authentication and user-management module, the foundation of the security and identity layer of the platform.
  \item \textbf{Chapter 4 -- Release 2 :} details the training catalog and the course-management module, which form the functional heart of the platform.
  \item \textbf{Chapter 5 -- Release 3 :} covers the reservation system, the notifications, the smart dashboard and the integrated chatbot.
  \item \textbf{Chapter 6 -- Realization :} presents the technical environment, the architecture, the development tools and the most representative interfaces of the final product.
\end{itemize}

\medskip
\noindent
A general conclusion closes this dissertation, summarizing the contributions of the work, evaluating its results against the initial objectives and tracing concrete perspectives of evolution for the platform.

\cleardoublepage
